\section{Glossary}
\label{sec:glossary}

Terms are grouped by where they come from. A term defined in the body carries
its section number.

\subsection*{Order theory}

\begin{description}[leftmargin=0pt,labelindent=0pt,style=nextline,
                    font=\normalfont\bfseries]
\item[Antisymmetry] If $x \leq y$ and $y \leq x$ then $x = y$. In a stored DAG
  this is exactly acyclicity (\S\ref{sec:orders}).
\item[Complete lattice] A poset in which \emph{every} subset --- not merely
  every pair --- has a meet and a join. Automatically has a top and a bottom.
\item[Cone] For a node $x$: the descendant cone $\cone{x}$ is everything
  reachable from $x$ (including $x$); the ancestor cone $\up{x}$ is everything
  from which $x$ is reachable. \odag{}'s query answers are intersections of
  descendant cones (\S\ref{sec:get}).
\item[Cover] $y$ covers $x$ when $x < y$ with nothing strictly between. The
  covering pairs are exactly the edges of a Hasse diagram, and exactly the
  edges of a transitive reduction.
\item[DAG] Directed acyclic graph. Reachability in a DAG is a partial order,
  and every finite poset arises this way (Prop.~\ref{prop:dag-poset}).
\item[Dedekind--MacNeille completion] The smallest complete lattice into which
  a poset embeds, preserving all meets and joins it already had. Equal to the
  concept lattice of $(P,P,\leq)$ (Thm~\ref{thm:dm}); \citet{macneille1937}.
\item[Galois connection] A pair of containment-reversing maps between the
  subsets of two sets, such that $A \subseteq g(B) \iff B \subseteq f(A)$.
  Both composites are closure operators (\S\ref{sec:closure}).
\item[Hasse diagram] A drawing of a poset showing only covering pairs, with
  greater elements higher. All other relationships are read by walking.
\item[Lattice] A poset in which every pair has a meet and a join.
\item[Meet / join] Greatest lower bound / least upper bound of a set of
  elements, when one exists.
\item[Partial order] Reflexive, antisymmetric, transitive relation. ``Partial''
  because two elements may be incomparable --- which is the point.
\item[Poset] A set with a partial order.
\item[Transitive closure] Add an edge wherever a path exists.
\item[Transitive reduction] Remove every edge implied by a path. For a DAG it
  is unique and equals the Hasse diagram (Thm~\ref{thm:unique-reduction};
  \citealp{aho1972}).
\end{description}

\subsection*{Formal Concept Analysis}

\begin{description}[leftmargin=0pt,labelindent=0pt,style=nextline,
                    font=\normalfont\bfseries]
\item[Attribute exploration] An interactive protocol in which the system asks
  an expert whether each unconfirmed implication always holds, and the expert
  either confirms it or supplies a counterexample
  (\S\ref{sec:learn-exploration}; \citealp{ganter2016}).
\item[Concept lattice $\Bcl(\ctx)$] All formal concepts of a context, ordered
  by extent containment. Always a complete lattice (Thm~\ref{thm:basic}).
\item[Contranominal scale] The context $(\{1..n\},\{1..n\},\neq)$. Every
  attribute set is closed, so it has $2^n$ concepts --- the canonical
  worst case (\S\ref{sec:algorithms}).
\item[Derivation ($'$)] $A'$ is the attributes shared by all of $A$; $B'$ is
  the objects having all of $B$. In \odag{} terms $B' = \cmd{get}(B)$
  (Prop.~\ref{prop:dictionary}).
\item[Duquenne--Guigues base (stem base)] The unique minimum-size set of
  implications generating all implications of a context
  \citep{guigues1986}.
\item[Extent / intent] The two halves of a formal concept: the objects it
  covers, and the attributes defining it.
\item[Formal concept] A pair $(A,B)$ with $A'=B$ and $B'=A$ --- a query in
  balance with its own answer (\S\ref{sec:dictionary}).
\item[Formal context $(G,M,I)$] Objects, attributes, and which has which. The
  same object as a binary code matrix (\S\ref{sec:code-is-context}).
\item[Iceberg lattice] The concepts above a support threshold
  \citep{stumme2002}.
\item[Implication $B \to C$] Every object with all of $B$ has all of $C$.
  Non-monotone: new data can destroy it (\S\ref{sec:learn-implications}).
\item[NextClosure] Ganter's algorithm enumerating all closed sets in lectic
  order, in constant memory \citep{ganter1984}.
\item[Pattern structure] \fca{} generalised so descriptions live in a
  meet-semilattice rather than being attribute sets; the interval case
  matches \odag{}'s typed values (\S\ref{sec:learn-patterns};
  \citealp{ganter2001,kaytoue2011}).
\item[Reduced labelling] Drawing convention where each object and attribute
  name appears once, at the concept that introduces it
  (\S\ref{sec:basic-theorem}).
\item[Scaling] Turning non-binary data into a context by choosing, per
  attribute, which comparisons matter (\S\ref{sec:learn-scaling}).
\item[Stability] The fraction of subsets of a concept's extent that still
  close to the same intent: robustness against losing objects
  \citep{kuznetsov2007}.
\end{description}

\subsection*{\odag{}}

\begin{description}[leftmargin=0pt,labelindent=0pt,style=nextline,
                    font=\normalfont\bfseries]
\item[Ancestor-set code] $\code{x} = \up{x}$, the reflexive ancestor set read
  as a bit vector. Equal to the \fca{} intent; determines the graph
  (Prop.~\ref{prop:roundtrip}).
\item[Canonical form] The unique transitive reduction plus
  denotation-canonicalised values plus a deterministic serialisation. What
  makes a single root hash possible (\S\ref{sec:canonical}).
\item[Certificate] A self-describing envelope carrying the raw record bytes
  needed to re-check an \cmd{is\_below} answer against a root alone
  (\S\ref{sec:agents-verify}).
\item[Cone summary] A published derived record stating a broad category's
  membership, so a thin client need not walk it.
\item[\texttt{descendant\_count}] The exact size of a node's cone,
  delta-maintained. \fca{}'s support, and the query planner's statistics.
\item[Dimension / parametric value] A typed value such as
  \texttt{weight(4kg)} whose order is computed from the name by exact
  rational arithmetic rather than stored as edges (\S\ref{sec:dimensions}).
\item[\texttt{get\_overlapping}] Candidate generation for parametric terms:
  complete for possibility, silent on satisfaction (G6).
\item[\texttt{is\_below}] Boolean subsumption, answered by an upward walk,
  fail-closed (G4).
\item[Merge] Union of assertions followed by re-reduction. Commutative and
  idempotent, hence a CRDT (Thm~\ref{thm:merge}).
\item[Pinned root / as-of] Asking a question against an immutable past state,
  which is how non-monotone questions become well defined (\S\ref{sec:asof}).
\item[Provenance record] A signed statement that an author asserted (or
  retracted, or endorsed) a claim, against the root they were looking at.
  Stored separately from the knowledge (\S\ref{sec:agents}).
\item[Residency] Whether a graph is eager (fully in memory), lazy (fetched as
  a query walks), or sparse (partially resident and writable). Semantics are
  identical across all three (G3).
\item[Root] The 32-byte content address of an entire knowledge state.
\item[\texttt{*}] The graph root node, implicit ancestor of everything.
\end{description}

\subsection*{Distributed systems and cryptography}

\begin{description}[leftmargin=0pt,labelindent=0pt,style=nextline,
                    font=\normalfont\bfseries]
\item[Content addressing] Retrieving a blob by the hash of its own bytes;
  nothing is ever overwritten.
\item[CRDT] Conflict-free replicated data type: replicas converge under a
  commutative, idempotent merge, with no coordination
  \citep{shapiro2011}.
\item[Feed] A signed, followable pointer to a changing value --- here, ``the
  latest root published by this key''.
\item[Merkle proof] The path of hashes showing a record is (or, given a
  canonical encoding, is not) part of a hashed structure
  \citep{merkle1987}.
\item[Monotone] An answer that can only grow as information arrives. Required
  of anything asked of a merging store (G2).
\item[Postage stamp] Prepaid storage rent on Swarm; the natural eviction
  mechanism for regenerable derived indexes (\S\ref{sec:crypto-economics}).
\end{description}
